\documentclass[12pt, titlepage]{article}
\usepackage[colorlinks=true, linkcolor=blue, urlcolor=magenta]{hyperref}

\title{Feedback Controls Lab Weekly Report \\ \normalsize{Week 10}}
\author{Aydan Vissing, Brady Schick, Gabriel Southwell, Simon Ross}
\date{\today}

\begin{document}
\maketitle

\section*{Aydan Vissing (9 hours spent)}

I spent a lot of time on the Mini-Project this week, along with a minor redesign of the frame that took longer than expected. The Mini-project i originally wanted to do ended up being too difficult for what little time I had so I switched part way through and finished the Voltage Regulator project instead. That took up about five hours. I spent another hour trying to figure out how to redesign the chassis so it would be easier to put the motors in and also have more strength in the parts that failed in the last prototype during flex testing and cleanup. The last three hours were spent in lab and our weekly meeting. I plan on finishing the python version of Dr. Fredette's mini-project sometime this week and getting the new frame printed since I was unable to check if the printer in the ENS was empty this weekend. 

\section*{Brady Schick (11 hours spent)}

I spent most of the week on continuing to implement the controller logic to the MATLAB simulation file from Dr. Fredette. It works at least to a point where some results can be given in the presentation. The rest of the time was spent finishing up the machine learning MP.

\section*{Gabriel Southwell (4 hours spent)}

This was a low hour week for me because of ACARCH and Theo and a few ofther things. Most weeks will not be like this.

This week I did the sensors mini project and worked a little more on the time of flight sensors. I talked to Dr. Kohl about them and we both think it would be best if I wrote my own library for them. I will look at existing esp32 libraries for inspiration on how it should work, but I am going to have to make a system to reassigning i2c addresses. This will require a unique, digital pin from the esp to every sensor though. We should have more than enough pins, but worst case scenario we could save a few with a binary decoder. 

This week I will work on that and work on our team progress presentation. 

\section*{Simon Ross (8 hours spent)}

This week has been the start of a big push on a couple important tasks (at least for me). The first is our upcoming presentation on Wednesday. This will basically be a status report in front of the entire class, so it will take some work. The second is the printed circuit board for our drone. I aim to have a draft of the PCB fully designed by the end of the next week, which I think is an attainable goal. This will require a lot of collaboration with Gabe regarding the pinout of the components. 

With regard to our schedule, we have made a slight change, working more on the PCB and frame than the prototype drone. This means that our upcoming November milestone is to have the PCB ordered instead of a functional prototype. We hope that the PCB will serve as a functional piece of equipment anyways.

\end{document}